\documentclass[10pt,a4paper]{amsart}
\usepackage[margin=19mm]{geometry}
\usepackage{amsmath,amssymb,mathtools}
\usepackage[scr=rsfs,cal=euler]{mathalfa}
\usepackage{xfrac}
\usepackage[unicode,hidelinks,bookmarks=false]{hyperref}
\newcommand{\ie}{i.e.\ }
\newcommand{\cf}{cf.\ }
\newcommand{\resp}{respectively\ }
\newcommand{\Subsection}{\subsection}
\providecommand{\nobreakdash}{\nobreak\hbox{-}}
\setlength{\emergencystretch}{2em}
\multlinegap0pt
\usepackage{bm}
\usepackage{bbm}
\usepackage{dsfont}
\usepackage{xspace}
\usepackage{cancel}
\usepackage{mathtools}
\usepackage{amsthm}
\makeatletter
\@ifundefined{theorem}{\newtheorem{theorem}{Theorem}}{}
\makeatother
\newcommand{\Spec}{\textnormal{Spec}}
\newcommand{\Z}{\mathbb{Z}}
\newcommand{\bbP}{\mathbf{P}}
\newcommand{\ca}[1]{{\mathcal{#1}}}
\newcommand{\pr}{\textnormal{pr}}
\title[A surface with representable $\text{CH}\_{0}$-group but no un]{A surface with representable $\text{CH}\_{0}$-group but no universal zero-cycle}
\author{Theodosis Alexandrou}
\date{Source: arXiv:2602.13435v2 (CC BY 4.0)}
\begin{document}
\maketitle

\noindent\emph{Source and licence.} A surface with representable $\text{CH}_{0}$-group but no universal zero-cycle. Version arXiv:2602.13435v2 (CC BY 4.0), distributed under the Creative Commons Attribution 4.0 International licence, as recorded in the arXiv licence metadata for this exact version and shown on the abstract page https://arxiv.org/abs/2602.13435v2. Full rights notice: licenses/arxiv-2602.13435-CC-BY-4.0.txt.\par\smallskip

\noindent\textit{Editorial scope.} Counted unit: the Theorem whose source label is \texttt{thm:degeneration\_S} in the article (source0.tex lines 1021--1027, source label \texttt{thm:degeneration\_S}), delivered together with the article's own complete original proof of it (source0.tex lines 1028--1075, one \texttt{proof} environment of 48 lines). No mathematics is written, repaired or re-derived: statement and proof are the article's own text line for line. Required context retained uncounted: none. Every definition, display and label of those lines is retained unchanged. Omitted from this package and not redistributed: 1--1020, 1076--1411. Nothing inside a retained excerpt is omitted or altered: every retained body is delivered exactly as the article's lines are written, so no omission receipt of any kind accompanies this package. Citations: the retained excerpts carry no citation command at all, so no bibliography entry of the article is transcribed and no reference is redistributed. Reference adaptation: none was needed and none was made; every reference inside the delivered unit resolves inside it, so no unresolved reference or citation is produced. Printed slips kept literal: no printed word, symbol, hypothesis or number of the counted statement or of its proof was altered. Layout and class: the article's own class is \texttt{amsart} with packages including mathtools, mathrsfs, fontenc, url, amsmath, array, graphicx, amsfonts, amssymb, amstext, amsthm, hyperref, enumerate, old-arrows; the wrapper is the lane's pinned \texttt{amsart} editorial wrapper at 10pt on A4, which restates the theorem-like environments the retained text needs and adds the article's own macro definitions that the retained text needs (\texttt{\textbackslash Spec}, \texttt{\textbackslash Z}, \texttt{\textbackslash bbP}, \texttt{\textbackslash ca}, \texttt{\textbackslash pr}), with extra wrapper packages (bm, bbm, dsfont, xspace, cancel, mathtools). The delivered numbering is the unit's own. Assets: no retained span contains a figure environment, a graphics-inclusion command, a PDF-page inclusion, a table or a TikZ picture; no image file is copied into this package, the package declares no asset dependency and no image file is redistributed with it. Licence: version arXiv:2602.13435v2 of this article is distributed under the Creative Commons Attribution 4.0 International licence, as recorded in the arXiv licence metadata for this exact version and shown on the abstract page https://arxiv.org/abs/2602.13435v2. A full rights notice is delivered at licenses/arxiv-2602.13435-CC-BY-4.0.txt. Characters: every retained excerpt is pure ASCII, so the delivered engine, XeLaTeX with its default Latin Modern text font, reports no missing character.
\begin{theorem}\label{thm:degeneration_S} Let $k$ be an algebraically closed field of characteristic different from $2,$ and let $E$ be an elliptic curve over $k$. Set $\Delta:=\Spec k[[t]],$ and fix an algebraic closure $K$ of its function field. Then there exists a regular, flat, projective scheme $\mathcal{S}\to\Delta$ together with a morphism $$\phi_{\mathcal{S}}\colon\mathcal{S}\longrightarrow E\times\Delta,$$ such that the following hold:
\begin{enumerate}
    \item\label{it:generic_fibre'} The geometric generic fibre $S_{K}$ is a bielliptic surface of type 2, and the induced morphism $\phi_{S_{K}}\colon S_{K}\to E\times_{k} K$ coincides with the Albanese fibration of $S_{K}$.
    \item\label{it:special_fibre'} The special fibre $S_{0}=R_{1}\cup R_{2}$ is a reduced simple normal crossings divisor on $\mathcal{S}$ with exactly two irreducible components. Each surface $R_{i}$ is a minimal ruled surface over an \'etale double cover $E_{i}\to E,$ and the two double covers are not isomorphic.\par
    The intersection $$C:=R_{1}\cap R_{2}$$ is a smooth, irreducible elliptic curve, embedded in each $R_{i}$ as a $2$-fold multisection. 
\end{enumerate}  
\end{theorem}

\begin{proof} Let $p_{\mathcal{F}}\colon\mathcal{F}\to\Delta$ be a minimal regular model of an elliptic curve with geometric generic fibre $F_{K}:=\mathcal{F}\times_{\Delta}K,$ and assume that the special fibre $F_{0}$ is of Kodaira type $\bf{I}_{4}$. Then $$F_{0}=C_{0}\cup C_{1}\cup C_{2}\cup C_{3},$$ where each $C_{i}\cong\bbP^{1}_{k}$. The components are arranged in a cycle: for $i\in\Z/4,$ $$C_{i}\cap C_{i+1}=\{\text{single node}\}, \quad C_{i}\cap C_{j}=\emptyset\ \text{if}\ j\not\equiv i\pm 1\ \text{(mod 4)}.$$
Each node $C_{i}\cap C_{i+1}$ is obtained by identifying the point $\infty\in C_{i}$ with the point $0\in C_{i+1}$.\par
The smooth locus $\mathcal{F}^{\mathrm{sm}}\to\Delta$ of $p_{\mathcal{F}}$ is the N\'eron model of its generic fibre. Let $\mathcal{F}^{\mathrm{sm}}[2]$ denote its $2$-torsion subgroup scheme; one has $$\mathcal{F}^{\mathrm{sm}}[2]\cong\mu_{2}\times(\Z/2)_{\Delta}.$$ This group scheme acts on $\mathcal{F},$ and the induced action on the special fibre is as follows: the factor $\mu_{2}$ acts on each component $C_{i}$ by multiplication by $\pm1,$ fixing the two nodes of $C_{i}$, while the factor $\Z/2$ permutes the components by sending $C_{i}$ to $C_{i+2}$.\par
Choose a non-trivial $2$-torsion class $$0\neq\epsilon\in\mathcal{F}^{\mathrm{sm}}[2]$$ whose restriction to the special fibre corresponds to the element $(1,\bar{1})\in\mu_{2}\times\Z/2,$ and let $\varphi_{\epsilon}\colon \ca F\longrightarrow\ca F$ denote translation by $\epsilon$.\par
The inversion morphism $$[-1]\colon F_{K}\longrightarrow F_{K},\quad x\longmapsto -x$$ extends canonically to an involution
\begin{equation*}
    \iota\colon \ca F\longrightarrow\ca F
\end{equation*}
over $\Delta$. On the special fibre $F_0,$ the automorphism $\iota$ induces, for each $i\in\Z/4,$ the isomorphism $$C_i\longrightarrow C_{-i},\quad (x_{0}:x_{1})\longmapsto (x_{1}:x_{0}).$$\par
Fix generators $a,b\in E[2],$ and let $\tau_{a},\tau_{b}\colon E\to E$ denote translation by $a$ and $b,$ respectively. Define
\begin{equation}\label{eq:inv_bielliptic}
    \sigma_{1}:=\tau_{a}\times_{k}\iota, \qquad 
    \sigma_{2}:=\tau_{b}\times_{k}\varphi_{\epsilon}.
\end{equation}
The involutions $\sigma_{i}$ act freely on $E\times_{k}\ca F,$ and they commute: $\sigma_{1}\circ\sigma_{2}=\sigma_{2}\circ\sigma_{1}$.\par
Set \begin{equation}\label{eq:deg_I}
\widetilde{\ca S}:=\bigl(E\times_{k}\ca F\bigr)/\sigma_{1}.\end{equation} The natural projection $\pr_{1}\colon E\times_{k}\ca F\to E\times_{k}\Delta$ descends to a morphism \begin{equation}\label{eq:albanese_I}\phi_{\widetilde{\ca S}}\colon \widetilde{\ca S}\longrightarrow \bigl(E/a\bigr)\times_{k}\Delta .\end{equation}
The family $\widetilde{\ca S}\to \Delta$ in \eqref{eq:deg_I} is a strictly semistable degeneration whose geometric generic fibre $\widetilde{S}_{K}$ is a bielliptic surface of type 1. The special fibre $\widetilde{S}_{0}$ has exactly three irreducible components $\widetilde{R}_{i},$ each a minimal elliptic ruled surface. More precisely, 
\begin{equation}\label{eq:components_I}
\widetilde{R}_{1} = \bigl(E \times C_0\bigr)/\sigma_1,\quad
\widetilde{R}_{2} =\bigl( E \times (C_1 \sqcup C_3)\bigr)/\sigma_1 \cong E \times \bbP^1_k,\quad
\widetilde{R}_{3} = \bigl(E \times C_2\bigr)/\sigma_1 .
\end{equation}

The dual graph of $\widetilde{S}_{0}$ is a chain with end components $\widetilde{R}_{1}$ and $ \widetilde{R}_{3}$. The intersections 
\begin{equation*}
    \widetilde{C}_{1,2}:= \widetilde{R}_{1}\cap \widetilde{R}_{2},\qquad \widetilde{C}_{2,3}:=\widetilde{R}_{2}\cap \widetilde{R}_{3}
\end{equation*}
are naturally isomorphic to the elliptic curve $E$. Moreover, the induced morphisms $$ \widetilde{C}_{1,2}\subset \widetilde{R}_{1} \longrightarrow E/a,\qquad \widetilde{C}_{2,3}\subset \widetilde{R}_{3} \longrightarrow E/a$$ coincide with the canonical double cover $$E\longrightarrow E/a.$$ In addition, the curves $\widetilde{C}_{i,i+1}$ embed as disjoint sections of the ruled surface $\widetilde{R}_{2}$.\par
The involution $\sigma_{2}$ in \eqref{eq:inv_bielliptic} descends to an automorphism $\widetilde{\sigma}_{2}$ of $\widetilde{\ca S}$. The quotient \begin{equation}\label{eq:deg_II}
    \ca S:=\widetilde{\ca S}/\widetilde{\sigma}_{2}\longrightarrow\Delta
\end{equation}
is a strictly semistable degeneration of a bielliptic surface of type 2. The involution $\widetilde{\sigma}_{2}$ identifies the components $\widetilde{R}_{1}$ and $\widetilde{R}_{3}$ and leaves $\widetilde{R}_{2}$ invariant. Under the isomorphism $\widetilde{R}_{2}\cong E \times \bbP^1_k$ in \eqref{eq:components_I}, its restriction is given by $$\sigma\colon\bigl(x,(x_{0}:x_{1})\bigr)\longmapsto \bigl(x+(a+b),(x_{1}:x_{0})\bigr).$$
It follows that the special fibre $S_{0}$ has exactly two irreducible components $R_{i},$ each a minimal elliptic ruled surface. More precisely, 
\begin{equation}\label{eq:components_II}
R_{1}\cong \widetilde{R}_{1}\cong \widetilde{R}_{3},\qquad R_{2}\cong \bigl(E\times\bbP^1_{k}\bigr)/\sigma.
\end{equation}
Their intersection $C:=R_{1}\cap R_{2}$ is isomorphic to $E,$ and embeds in each ruled surface $R_{i}$ as a double cover of its base $E_{i}$. In the notation of \eqref{it:special_fibre'}, the corresponding \'etale double covers $E_{i}\to E,\ (i=1,2),$ are given by \begin{equation}\label{eq:double_covers}
    E/a\longrightarrow E/E[2]\cong E,\qquad E/(a+b)\longrightarrow E/E[2]\cong E,
\end{equation}
respectively.\par
Finally, the morphism $\phi_{\widetilde{\ca S}}$ defined in \eqref{eq:albanese_I} descends, upon passage to the quotient by $\widetilde{\sigma}_{2}$, to a morphism 
\begin{equation}\label{eq:albanese_II}
\phi_{\ca S}\colon \ca S\longrightarrow \bigl(E/E[2]\bigr)\times_{k}\Delta\cong E\times_{k}\Delta,
\end{equation}
which restricts on the geometric generic fibre to the Albanese fibration of the bielliptic surface $S_{K}$.\par
This completes the proof.
\end{proof}

\end{document}
